\begin{tikzpicture}[x=4.25cm,y=1.0cm]
\draw[-{Latex[length=1.4mm]}] (0,0) -- (0,3.15) node[above,font=\scriptsize] {flux [arb.]};
\draw[-{Latex[length=1.4mm]}] (-0.01,0) -- (1.10,0) node[below,font=\scriptsize] {$\xi$};
\foreach \xx in {0.5,1.0}{\draw (\xx,0.03) -- (\xx,-0.03) node[below,font=\scriptsize] {\xx};}
\draw[black!45] (0,1.0) -- (1,0);
\draw[black!45,densely dashed] (0,0) -- (1,2.0);
\fill[black!45] (0.3333,0.6667) circle (0.75pt);
\node[font=\scriptsize,black!55] at (0.22,2.35) {$A<0$};
\draw[black!65] (0,1.5) -- (1,0);
\draw[black!65,densely dashed] (0,0) -- (1,1.5);
\fill[black!65] (0.5000,0.7500) circle (0.75pt);
\node[font=\scriptsize,black!65] at (0.50,1.55) {$A=0$};
\draw[thick] (0,2.0) -- (1,0);
\draw[thick,densely dashed] (0,0) -- (1,1.0);
\draw[densely dotted] (0.6667,0) -- (0.6667,0.6667);
\fill (0.6667,0.6667) circle (0.95pt);
\node[below,font=\scriptsize] at (0.6667,-0.15) {$\xi_\eq=2/3$};
\node[right,font=\scriptsize] at (0.04,2.08) {forward $r^+(1-\xi)$};
\node[right,font=\scriptsize] at (0.65,0.50) {reverse $r^-\xi$};
\draw[-{Latex[length=1.3mm]},semithick] (0.39,0.98) -- (0.59,0.76);
\node[font=\scriptsize,align=center] at (0.78,2.45) {$A$ raises $r^+/r^-$\\intersection moves right};
\end{tikzpicture}
\caption{정지점 유도. 정방향 플럭스 $r^+(1-\xi)$(감소 직선)와 역방향 $r^-\xi$(증가 직선)의 교점이 정지점 $\xi_\eq$(식~\eqref{eq:logisticsolve})이다. 기존 두 직선 구도에 $A<0$, $A=0$, $A>0$ 세 케이스를 겹쳐 affinity가 $r^+/r^-$를 바꾸며 교점을 이동시키는 점을 보강했다. 검정 케이스는 $\mathcal A>0$, $r^+=2r^-$라 $\xi_\eq=2/3$이고, 회색 케이스는 detailed balance(식~\eqref{eq:db})의 기준이다. 좌표는 \code{T5\_calc.py}가 $\xi_\eq=r^+/(r^++r^-)$로 계산했다.}
\label{fig:flux}
